% s-078-p078-m003; original PDF 78, 79
\Aphorism{s-078-p078-m003}{70}
Hang zur Einsamkeit kann zweierlei sein: Jünglingssehnsucht oder Manneswille, „Flucht des Kranken oder Flucht vor Kranken“, geistige Pubertät oder geistige Reife. Nur die eine Art hält es in der wirklichen Einsamkeit auf die Dauer aus; die andere sucht nur eine Theater-Einsamkeit, in der man gesehen wird, und taucht daher immer wieder in den grossen Markt- und Verkehrsstrudeln auf, um sich der nöthigen Zuschauer zu versichern. Diese unechte Art überwiegt, und im Ganzen ist Einsamkeit das bestmissbrauchte Glücksgeschenk und Privilegium des Menschenvolkes — allenfalls die \Sperr{Musik} ausgenommen, die nicht minder selten an den rechten Mann kommt und die auch ihrerseits in zwei Arten zerfällt, eine für Männer und eine für —; Musik als \Foreign{virtus} und Musik als \Foreign{voluptas.}
\par

